\chapter{Change points and process control}\label{ch:change-points}

Operations ask of a daily series questions of different shapes: was yesterday normal, did the level of activity change
at some date, and should someone be woken up this hour? The exploratory notebooks answer them with robust day rules,
penalised change-point detection, control charts and hourly limits, and on this dataset every answer is the same: nothing
happened. A null result is only as informative as the method's ability to see what it is looking for. This chapter
therefore pairs every null finding with the size of effect the method would have caught, and checks that each monitor
alerts at the rate its design promises when nothing happens.

All series here are counted on each process's own delivery clock (\cref{ch:time-and-clocks}), per customer country,
over the $\est{ch09.days}{value}$ complete delivery days.

\section{The daily residual}\label{sec:cp-residual}

Let $y_t$ be the number of transactions of a market on delivery day $t$. Its logarithm is modelled as
\begin{equation}\label{eq:cp-calendar}
  \log y_t = \alpha_{w(t)} + \beta\, \mathbb{1}\{t \text{ ends a month}\} + \gamma\, t + e_t ,
\end{equation}
with a coefficient per weekday $w(t)$ and a linear trend. \Cref{ch:time-and-clocks} found that, after the weekday,
holidays and pay cycles explain nothing that survives the control of the false-discovery rate, so this smaller model
loses nothing material. The residuals $e_t$ have a standard deviation between $\est{ch09.resid-sd-min}{value}$ and
$\est{ch09.resid-sd-max}{value}$ on the log scale, about a tenth of the day's volume, and their lag-one autocorrelation
is at most $\est{ch09.lag1-max}{value}$ in absolute value: once the week is removed, consecutive days are independent.
That property is what makes the monitors below calibratable, and it is itself a fingerprint of a generator, since
real daily volumes carry persistent shocks.

\paragraph{Abnormal days.} The country notebooks call a day abnormal when its residual lies more than
$\est{ch09.z-day}{value}$ robust standard deviations ($\kappa\,\mathrm{MAD}$, \cref{eq:sf-scale}) from the median. On the residuals of
\cref{eq:cp-calendar} no day in any market qualifies ($\est{ch09.abnormal}{value}$), against
$\est{ch09.abnormal-expected}{value}$ expected by chance under normality; the notebooks' fuller calendar model also finds
$\est{ch09.nb-abnormal}{value}$. There are no outages, campaign peaks or delivery failures in the data.

\section{Level shifts}\label{sec:cp-shifts}

\subsection{Penalised segmentation}

A level shift is a change in the mean of a series. For a series $x_1, \dots, x_n$ (here the weekly means of $e_t$,
$\est{ch09.weeks}{value}$ weeks per market) and change points $0 = \tau_0 < \tau_1 < \dots < \tau_K < \tau_{K+1} = n$,
penalised segmentation solves
\begin{equation}\label{eq:cp-penalised}
  \min_{K,\,\tau} \; \sum_{k=0}^{K} C\bigl(x_{(\tau_k + 1) : \tau_{k+1}}\bigr) + \beta K ,
\end{equation}
where $C$ is the cost of a segment and $\beta$ the price of a change. PELT finds the exact minimum in time close to
linear by discarding candidate change points that can never be optimal again \autocite{killick2012}. The method has
inputs that decide what it can find: the cost and the penalty. With the Gaussian cost $C(x_{a:b}) = \sum_{i=a}^{b} (x_i - \bar x_{a:b})^2$ the
penalty is in units of variance, and $\beta = 2 \hat\sigma^2 \log n$ is the Schwarz criterion for a mean shift in
Gaussian noise \autocite{schwarz1978}. The country notebooks instead use a kernel cost with a Gaussian (RBF) kernel
\autocite{arlot2019}, which can detect changes in the whole distribution, with a penalty of $3 \log n$. Each term of a
kernel cost is bounded, and a small shift of the mean barely changes the kernel similarities, so in kernel units this
penalty is a high bar for a mean shift; it is not calibrated to the noise of the series. The last detector is the
permutation binary segmentation of \cref{sec:sf-detectors}, at level $\est{ch09.bs-alpha}{value}$, which is valid here
because the weekly residuals are close to independent.

\subsection{What a null result is worth}

On the real series none of the detectors finds a change point: the kernel PELT of the notebooks
($\est{ch09.found-rbf}{value}$, as the notebooks report, $\est{ch09.nb-change-points}{value}$), the Gaussian PELT
($\est{ch09.found-l2}{value}$) and the binary segmentation ($\est{ch09.found-bs}{value}$). To learn what this rules
out, steps of known size were planted at random weeks of the real residual series, $\est{ch09.sims}{value}$ times per size
and market, and a step counted as found when a change point fell within $\est{ch09.tol}{value}$ weeks of it; the
false-alarm rate was measured on random permutations of the same series, which keep their values and destroy any
change (\cref{tab:ch09-shift-detectors,fig:ch09-power}).

\input{../build/generated/tables/ch09-shift-detectors}

\begin{figure}[htbp]
  \centering
  \includegraphics[width=0.8\textwidth]{fig-09-power}
  \caption{Share of planted level steps found by each detector, by the size of the step in weekly log volume, pooled
  over the markets. The dotted line marks the conventional target power of
  $\est{ch09.power-target}{value}\,\%$.}
  \label{fig:ch09-power}
\end{figure}

The kernel PELT never raises a false alarm ($\est{ch09.false-rbf}{pct}\,\%$, interval $\est{ch09.false-rbf}{pct_ci}$), but
it is also blind to moderate steps: a step of $\est{ch09.step-small}{value}\,\%$ is found
$\est{ch09.power-rbf-5}{pct}\,\%$ of the time, a step of $\est{ch09.step-mid}{value}\,\%$ $\est{ch09.power-rbf-10}{pct}\,\%$
($\est{ch09.power-rbf-10}{pct_ci}$), and only steps of $\est{ch09.mde-rbf}{value}\,\%$ and more reach the target power. The Gaussian PELT, with a false-alarm rate of $\est{ch09.false-l2}{pct}\,\%$ ($\est{ch09.false-l2}{pct_ci}$), and the
binary segmentation, with $\est{ch09.false-bs}{pct}\,\%$ ($\est{ch09.false-bs}{pct_ci}$), both find steps of
$\est{ch09.mde-l2}{value}\,\%$ at the target power ($\est{ch09.power-l2-5}{pct}\,\%$ and $\est{ch09.power-bs-5}{pct}\,\%$).
The null result of the notebooks therefore stands, and it is stronger than they could show: the level of transactions did
not shift by more than about $\est{ch09.mde-l2}{value}\,\%$ in any market at any time, roughly the standard deviation of a weekly mean
($\est{ch09.weekly-sd-max}{value}$ on the log scale). The comparison also settles a design question: for monitoring the
level of a series, a cost matched to the change sought with a penalty in the units of its noise beats a general-purpose
cost with an uncalibrated penalty, which is blind to the moderate shifts that matter most.

\section{Control charts}\label{sec:cp-charts}

A change-point method looks back; a control chart decides each day. The granularity notebooks chart the residuals of
daily volume with an exponentially weighted moving average \autocite{roberts1959},
\begin{equation}\label{eq:cp-ewma}
  z_t = \lambda e_t + (1 - \lambda) z_{t-1}, \qquad
  \operatorname{Var}(z_t) \to \sigma^2 \frac{\lambda}{2 - \lambda},
\end{equation}
with $\lambda = \est{ch09.lambda}{value}$ and limits at $\pm \est{ch09.ewma-l}{value}$ stationary standard deviations.
If the residuals are independent and normal, the share of days beyond the limits is
$2\Phi(-L) = \est{ch09.nominal-alert-pct}{value}\,\%$; alerts cluster in runs, so the average time between first alerts is
longer \autocite{lucas1990}. Both statements need independent residuals, which \cref{sec:cp-residual} established (the
charted series have lag-one autocorrelations of at most $\est{ch09.ewma-lag1-max}{value}$). Over
$\est{ch09.ewma-days}{value}$ market-days of transactions and contacts the charts alert on $\est{ch09.ewma-alerts}{value}$
days, against $\est{ch09.ewma-expected}{value}$ at the nominal rate (\cref{tab:ch09-charts}).

The daily decline rate is charted with a $p$-chart, whose limits $\bar p \pm 3\sqrt{\bar p (1 - \bar p) / n_t}$ widen on
days with fewer transactions \autocite{montgomery2012}. They are right only if the declines of a day are binomial
around a common rate; otherwise every day's extra-binomial variation becomes an alert, and the remedy is to inflate the
limits by the observed dispersion \autocite{laney2002}. The Pearson dispersion of the decline rate
(\cref{eq:sf-dispersion}) is at most $\est{ch09.dispersion-max}{value}$, so the binomial limits are calibrated as they
stand. They alert on $\est{ch09.p-alerts}{value}$ of $\est{ch09.p-days}{value}$ market-days against
$\est{ch09.p-expected}{value}$ expected; a count this high has a probability of $\est{ch09.p-tail}{value}$ under the nominal
rate, weak evidence of a slight excess that the near-unit dispersion does not explain and that one alert every few
months per market would not make troublesome. The contrast with the
campaign empty shares of \cref{sec:sf-empty}, whose dispersion was orders of magnitude larger over all rows, is the
practical lesson: measure the dispersion before trusting binomial limits.

\input{../build/generated/tables/ch09-charts}

\section{Hourly limits}\label{sec:cp-hours}

A streaming monitor judges each hour's count against limits that a correct model would breach a thousandth of the time
($\est{ch09.hour-nominal-pct}{value}\,\%$). With a Poisson model for the hour the limits are too narrow, and the reason is
structural.

\begin{proposition}\label{prop:cp-nb}
Let a day's total $N$ have mean $\mu$ and variance $\mu + c\mu^2$, and let each of its events fall in a given hour with
probability $p$, independently. Then the hour's count $X$ has mean $m = p\mu$ and variance $m + c\, m^2$.
\end{proposition}

\begin{proof}
Given $N$, $X$ is binomial with $\mathbb{E}(X \mid N) = pN$ and $\operatorname{Var}(X \mid N) = p(1-p)N$. By the law of
total variance, $\operatorname{Var}(X) = p(1-p)\mu + p^2(\mu + c\mu^2) = p\mu + c(p\mu)^2$.
\end{proof}

A day-level variation beyond Poisson therefore reaches every hour, with the same coefficient $c$. If the day's level is
gamma-distributed around its mean, the hour is negative binomial with size $1/c$, the classical gamma--Poisson mixture
\autocite{greenwood1920}. Both premises hold here. Hours are a multinomial thinning of their day: the dispersion of
hour counts given the day's total is $\est{ch09.hour-phi-ar}{value}$, $\est{ch09.hour-phi-co}{value}$ and
$\est{ch09.hour-phi-mx}{value}$ in the markets. The day's level varies beyond Poisson with a coefficient of variation
of $\est{ch09.cv-mx}{value}\,\%$ in Mexico and nearly the same elsewhere, the residual spread of
\cref{sec:cp-residual} seen from the other side. The variance of an hour is then $1 + c\,m$ times its Poisson
variance, so Poisson limits fail worst where hours are busiest (\cref{tab:ch09-hours}): they flag
$\est{ch09.pois-ar}{pct}\,\%$ of hours in Argentina, $\est{ch09.pois-co}{pct}\,\%$ in Colombia and
$\est{ch09.pois-mx}{pct}\,\%$ in Mexico ($\est{ch09.pois-mx}{pct_ci}$). Negative-binomial limits with the day-level $c$
flag $\est{ch09.nb-ar}{pct}\,\%$, $\est{ch09.nb-co}{pct}\,\%$ and $\est{ch09.nb-mx}{pct}\,\%$, below the nominal rate:
counts are integers, so a quantile limit can only be placed between them and the realised tail probability falls below
the nominal one. A monitor slightly conservative by construction is the right side to err on.

\input{../build/generated/tables/ch09-hours}

\section{Verdict and implications}\label{sec:cp-verdict}

Nothing happened in the activity of this bank: no abnormal day, no level shift larger than about $\est{ch09.mde-l2}{value}\,\%$, no
special cause in the control charts. The methods used to establish this are sound once calibrated, and the calibration
is the transferable result. Change-point detection for monitoring should use a cost matched to the change sought and a
penalty in the units of the noise, and should report the smallest shift it would have found. Control charts on counts
and rates need independent residuals and a dispersion check before binomial or Poisson limits are trusted. Hourly
limits must carry the day-level dispersion, as \cref{prop:cp-nb} shows, which turns a monitor that cries wolf several
times a day into one that is quiet when nothing happens. On real data these monitors are where operations would first see
an outage, a fraud wave or a missed delivery; on this dataset they can only be shown to be calibrated, not useful.

\section{Limitations}\label{sec:cp-limits}

The residual model omits holidays and pay cycles, which \cref{ch:time-and-clocks} found immaterial; the notebooks' fuller
calendar model gives the same null results. Planted steps are permanent shifts of the mean; a transient episode, a
change in variance or a gradual drift would need its own power analysis. The power curves pool the markets and use one
tolerance window, so a detector that places its change point a few weeks late is counted as missing it. The control-chart
and hourly results are recomputed here on transactions and contacts and stand in for the notebooks' tables, which were
not stored.
